%====================== FRONT MATTER ======================
\chapter*{Remerciements}
\addcontentsline{toc}{chapter}{Remerciements}
\markboth{Remerciements}{Remerciements}

Au terme de ce projet de fin d'études, je tiens à exprimer ma sincère gratitude
à mon encadrant pédagogique [\textit{à compléter}] pour son suivi rigoureux et
ses conseils déterminants, ainsi qu'à mon encadrant professionnel
[\textit{à compléter}] pour son expertise et sa disponibilité. Je remercie
l'ensemble du corps professoral de l'ENSIAS pour la qualité de la formation
d'ingénieur, et les membres du jury qui ont accepté d'évaluer ce travail. Je
dédie enfin ce travail à ma famille, dont le soutien a rendu ce parcours possible.

%====================== RÉSUMÉ ======================
\chapter*{Résumé}
\addcontentsline{toc}{chapter}{Résumé}
\markboth{Résumé}{Résumé}

Les équipes d'ingénierie logicielle accumulent une connaissance technique
volumineuse et fragmentée — spécifications, documents d'architecture, comptes
rendus et directives internes. La recherche par mots-clés en capture mal
l'intention sémantique, tandis que les assistants génératifs généralistes
produisent des réponses fluides mais \emph{non fondées} : hallucinations, absence
de citations et opacité du raisonnement.

\devpilot{} est une plateforme agentique de \emph{Génération Augmentée par
Récupération} (RAG) de qualité production qui répond à ces limites. Chaque
réponse est ancrée dans les documents propres à l'organisation, étayée par des
citations, évaluée automatiquement (fidélité, pertinence, risque d'hallucination)
et entièrement traçable. La génération repose sur un flux agentique de sept
agents coopérants, chacun persisté comme une exécution inspectable.

Au-delà du moteur RAG, le projet apporte trois contributions distinctives : une
tour de contrôle d'exploitation (\textbf{RAGOps}) pour l'observabilité du système
d'IA, un sous-système de \textbf{traçabilité} reconstruisant le cycle de vie
complet de chaque requête, et un \textbf{AI Execution Studio} qui visualise le
pipeline en temps réel. La plateforme est bâtie sur FastAPI, PostgreSQL, Redis,
Celery, Qdrant, Ollama et Gemini côté serveur, et Next.js côté client, le tout
orchestré par Docker Compose. Ce mémoire détaille la démarche d'ingénierie, la
conception, l'architecture et la validation du système.

\noindent\textbf{Mots-clés :} RAG agentique, grands modèles de langage, recherche
vectorielle, observabilité de l'IA, traçabilité, évaluation, FastAPI, Qdrant,
Next.js.

%====================== ABSTRACT ======================
\chapter*{Abstract}
\addcontentsline{toc}{chapter}{Abstract}
\markboth{Abstract}{Abstract}

Software engineering teams accumulate large, fragmented technical knowledge.
Keyword search misses semantic intent, while general-purpose generative
assistants produce fluent but \emph{ungrounded} answers — hallucinations, no
citations, no visibility into the reasoning.

\devpilot{} is a production-grade \emph{agentic Retrieval-Augmented Generation}
(RAG) platform that closes this gap. Every answer is grounded in the
organization's own documents, supported by citations, automatically evaluated
(faithfulness, relevance, hallucination risk) and fully traceable. Generation is
an agentic workflow of seven cooperating agents, each persisted as an inspectable
run.

Beyond the RAG engine, the project contributes three distinctive elements: an
operational control tower (\textbf{RAGOps}) for AI-system observability, a
\textbf{traceability} subsystem that reconstructs each query's full life cycle,
and an \textbf{AI Execution Studio} that visualizes the pipeline in real time.
The platform is built on FastAPI, PostgreSQL, Redis, Celery, Qdrant, Ollama and
Gemini on the server, with Next.js on the client, orchestrated by Docker Compose.
This report details the engineering process, design, architecture and validation.

\noindent\textbf{Keywords:} Agentic RAG, Large Language Models, Vector Search, AI
Observability, Traceability, Evaluation, FastAPI, Qdrant, Next.js.

%====================== INTRODUCTION GÉNÉRALE ======================
\chapter*{Introduction générale}
\addcontentsline{toc}{chapter}{Introduction générale}
\markboth{Introduction générale}{Introduction générale}

\section*{Contexte}
L'intelligence artificielle générative a transformé l'accès à l'information, mais
son adoption en entreprise se heurte à une exigence non satisfaite par les
modèles généralistes : répondre à partir du patrimoine documentaire
\emph{propre} à l'organisation, avec des garanties de véracité et de
transparence. Pour une équipe d'ingénierie, une réponse plausible mais fausse
constitue un risque opérationnel, non une simple imperfection.

\section*{Motivation}
La promesse du RAG est de réconcilier la puissance générative des grands modèles
de langage avec la connaissance vérifiable de l'organisation~\cite{lewis2020}.
Transformer cette promesse en \emph{système de production} suppose toutefois de
résoudre des problèmes d'ingénierie concrets : ingestion asynchrone, recherche
vectorielle, ancrage et citation des sources, mesure de la qualité et, surtout,
observabilité et traçabilité de bout en bout.

\section*{Problématique}
Comment concevoir et réaliser une plateforme d'IA capable de centraliser les
connaissances d'une équipe, d'en récupérer l'information pertinente, de générer
des réponses fondées et citées, d'en mesurer la qualité, et d'offrir une
traçabilité et une observabilité permettant d'auditer et d'exploiter le système
comme tout service critique d'entreprise~?

\section*{Objectifs}
Le projet vise à concevoir et implémenter \devpilot{}, plateforme RAG agentique
multi-espaces offrant : espaces de travail isolés, ingestion documentaire
asynchrone, récupération hybride avec reclassement, génération agentique ancrée
et citée, évaluation automatique, \emph{streaming} temps réel, et un sous-système
d'observabilité (RAGOps, explorateur et inspecteur de traces, AI Execution
Studio).

\section*{Méthodologie}
La démarche est itérative et orientée production : elle s'est déroulée en
\textbf{vingt-trois phases} incrémentales, de la fondation conteneurisée jusqu'aux
fonctionnalités phares de visualisation. Chaque phase fixe un objectif vérifiable,
introduit des fonctionnalités, justifie ses choix techniques et est validée par
des tests. Le chapitre~\ref{chap:archi} (architecture et réalisation) constitue
le c{\oe}ur du travail.

\section*{Organisation du rapport}
Le chapitre~\ref{chap:contexte} pose le contexte, l'analyse de l'existant et le
cahier des charges. Le chapitre~\ref{chap:conception} présente l'analyse et la
conception (UML, modèle de données). Le chapitre~\ref{chap:archi} détaille
l'architecture et la réalisation. Le chapitre~\ref{chap:techno} justifie les
technologies et l'implémentation. Le chapitre~\ref{chap:validation} traite de la
validation et des résultats. Le rapport se clôt par une conclusion générale et
des perspectives.
